\subsection{Base Case}


% ----------------------------------- Sample RMSE by Horizon
\begin{figure}[htbp]
\centering
\caption{Sample RMSE by Horizon}
\label{fig:BL_SAMPLE_RMSE}

\begin{subfigure}[t]{0.5\textwidth}
    \centering
    \caption{\(\rho = 0.5\)}
    \includegraphics[width=\textwidth]{BASE_RMSE_BY_HORIZON_rho=0.5.png}
    \label{fig:BL_SAMPLE_RMSE_rho05}
\end{subfigure}

\begin{subfigure}[t]{0.5\textwidth}
    \centering
    \caption{\(\rho = 0.7\)}
    \includegraphics[width=\textwidth]{BASE_RMSE_BY_HORIZON_rho=0.7.png}
    \label{fig:BL_SAMPLE_RMSE_rho07}
\end{subfigure}

\begin{subfigure}[t]{0.5\textwidth}
    \centering
    \caption{\(\rho = 0.95\)}
    \includegraphics[width=\textwidth]{BASE_RMSE_BY_HORIZON_rho=0.95.png}
    \label{fig:BL_SAMPLE_RMSE_rho095}
\end{subfigure}

\begin{minipage}{0.9\textwidth}
\small
\textit{Note.} Each panel reports RMSE by horizon for the base DGP with varying $T$. Sub-graphs correspond to persistence values $\rho \in \{0.5,0.7,0.95\}$.
\end{minipage}

\end{figure}




\subsection{Extension 1}






% ----------------------------------- Small Sample Deviation from Truth
\begin{figure}[htbp]
\centering
\caption{Small-Sample Deviation from the True IRF}
\label{fig:DYN_MISSPEC_SMALL_SAMPLE_DEV}

\begin{subfigure}[t]{0.6\textwidth}
    \centering
    \caption{\(\rho = 0.5\)}
    \includegraphics[width=\textwidth]{DYN_MISSPEC_LIMITED_DEVIATION_rho=0.5_T=250_B=5000.png}
    \label{fig:DYN_MISSPEC_SMALL_SAMPLE_DEV_rho05}
\end{subfigure}

\begin{subfigure}[t]{0.6\textwidth}
    \centering
    \caption{\(\rho = 0.7\)}
    \includegraphics[width=\textwidth]{DYN_MISSPEC_LIMITED_DEVIATION_rho=0.7_T=250_B=5000.png}
    \label{fig:DYN_MISSPEC_SMALL_SAMPLE_DEV_rho07}
\end{subfigure}

\begin{subfigure}[t]{0.6\textwidth}
    \centering
    \caption{\(\rho = 0.95\)}
    \includegraphics[width=\textwidth]{DYN_MISSPEC_LIMITED_DEVIATION_rho=0.95_T=250_B=5000.png}
    \label{fig:DYN_MISSPEC_SMALL_SAMPLE_DEV_rho095}
\end{subfigure}

\begin{minipage}{0.9\textwidth}
\small
\textit{Note.} Deviation from truth, \(\hat\theta_h - \theta_h\), at \(T=250\) over \(B=5{,}000\) replications, for VAR(4) (black) and LP(4) (crimson). Solid line is the mean deviation (finite-sample bias); shaded area is the central \(90\%\) Monte Carlo band (empirical 5th--95th percentile), not an SE band. Within each panel, sub-axes correspond to \(\tau \in \{0,0.1,0.3,0.6\}\); panels to persistence \(\rho \in \{0.5,0.7,0.95\}\).
\end{minipage}

\end{figure}


% ----------------------------------- Asymptotic Deviation from Truth
\begin{figure}[htbp]
\centering
\caption{Asymptotic Deviation from the True IRF}
\label{fig:DYN_MISSPEC_ASYMPTOTIC}

\begin{subfigure}[t]{0.6\textwidth}
    \centering
    \caption{\(\rho = 0.5\)}
    \includegraphics[width=\textwidth]{DYN_MISSPEC_ASYMPTOTIC_DEVIATION_rho=0.5_T=10000_B=1000.png}
    \label{fig:DYN_MISSPEC_ASYMPTOTIC_rho05}
\end{subfigure}


\begin{subfigure}[t]{0.6\textwidth}
    \centering
    \caption{\(\rho = 0.7\)}
    \includegraphics[width=\textwidth]{DYN_MISSPEC_ASYMPTOTIC_DEVIATION_rho=0.7_T=10000_B=1000.png}
    \label{fig:DYN_MISSPEC_ASYMPTOTIC_rho07}
\end{subfigure}

\begin{subfigure}[t]{0.6\textwidth}
    \centering
    \caption{\(\rho = 0.95\)}
    \includegraphics[width=\textwidth]{DYN_MISSPEC_ASYMPTOTIC_DEVIATION_rho=0.95_T=10000_B=1000.png}
    \label{fig:DYN_MISSPEC_ASYMPTOTIC_rho095}
\end{subfigure}

\begin{minipage}{0.9\textwidth}
\small
\textit{Note.} As above, at \(T=10{,}000\) over \(B=1{,}000\) replications. LP(4) (crimson) lies on the zero line at all horizons, whereas VAR(4) (black) retains a small, non-vanishing short-horizon bias. Solid line is the mean deviation; shaded area is the central \(90\%\) Monte Carlo band. Sub-axes within each panel correspond to \(\tau \in \{0,0.1,0.3,0.6\}\); panels to \(\rho \in \{0.5,0.7,0.95\}\).
\end{minipage}

\end{figure}


% ----------------------------------- Sample RMSE by Horizon X
\begin{figure}[htbp]
\centering
\caption{Sample RMSE by Horizon for Dynamic Misspecification}
\label{fig:DYN_MISSPEC_SAMPLE_RMSE}

\begin{subfigure}[t]{0.5\textwidth}
    \centering
    \caption{\(\rho = 0.5\)}
    \includegraphics[width=\textwidth]{DYN_MISSPEC_RMSE_BY_HORIZON_rho=0.5_T=250.png}
    \label{fig:DYN_MISSPEC_SAMPLE_RMSE_rho05}
\end{subfigure}

\begin{subfigure}[t]{0.5\textwidth}
    \centering
    \caption{\(\rho = 0.7\)}
    \includegraphics[width=\textwidth]{DYN_MISSPEC_RMSE_BY_HORIZON_rho=0.7_T=250.png}
    \label{fig:DYN_MISSPEC_SAMPLE_RMSE_rho07}
\end{subfigure}

\begin{subfigure}[t]{0.5\textwidth}
    \centering
    \caption{\(\rho = 0.95\)}
    \includegraphics[width=\textwidth]{DYN_MISSPEC_RMSE_BY_HORIZON_rho=0.95_T=250.png}
    \label{fig:DYN_MISSPEC_SAMPLE_RMSE_rho095}
\end{subfigure}

\begin{minipage}{0.9\textwidth}
\small
\textit{Note.} Each panel reports RMSE by horizon for the dynamically misspecified DGP with varying $T$. Sub-graphs correspond to persistence values $\rho \in \{0.5,0.7,0.95\}$.
\end{minipage}

\end{figure}





\subsection{Extension 2}
% ----------------------------------- Sample RMSE by Horizon X
\begin{figure}[htbp]
\centering
\caption{Sample RMSE by Horizon under Nonlinear Lag Misspecification}
\label{fig:NON_LIN_SAMPLE_RMSE}

\begin{subfigure}[t]{0.9\textwidth}
    \centering
    \caption{\(\rho = 0.5\)}
    \includegraphics[width=\textwidth]{NON_LIN_RMSE_BY_HORIZON_rho=0.5_T=250.png}
    \label{fig:NON_LIN_SAMPLE_RMSE_rho05}
\end{subfigure}

\begin{subfigure}[t]{0.9\textwidth}
    \centering
    \caption{\(\rho = 0.7\)}
    \includegraphics[width=\textwidth]{NON_LIN_RMSE_BY_HORIZON_rho=0.7_T=250.png}
    \label{fig:NON_LIN_SAMPLE_RMSE_rho07}
\end{subfigure}

\begin{minipage}{0.9\textwidth}
\small
\textit{Note.} Each panel reports RMSE by horizon for the nonlinear lag-quadratic DGP with $T=250$. Rows correspond to persistence values $\rho \in \{0.5,0.7\}$. Columns within each panel correspond to nonlinearity strength $\gamma$. Estimator definitions and performance measures are identical to the baseline design.
\end{minipage}

\end{figure}

 % ----------------------------------- Asymptotic RMSE by Horizon X
\begin{figure}[htbp]
\centering
\caption{Asymptotic RMSE by Horizon under Nonlinear Lag Misspecification}
\label{fig:NON_LIN_ASYMPTOTIC_RMSE}

\begin{subfigure}[t]{0.9\textwidth}
    \centering
    \caption{\(\rho = 0.5\)}
    \includegraphics[width=\textwidth]{NON_LIN_RMSE_BY_HORIZON_rho=0.5_T=10000.png}
    \label{fig:NON_LIN_ASYMPTOTIC_RMSE_rho05}
\end{subfigure}

\begin{subfigure}[t]{0.9\textwidth}
    \centering
    \caption{\(\rho = 0.7\)}
    \includegraphics[width=\textwidth]{NON_LIN_RMSE_BY_HORIZON_rho=0.7_T=10000.png}
    \label{fig:NON_LIN_ASYMPTOTIC_RMSE_rho07}
\end{subfigure}

\begin{minipage}{0.9\textwidth}
\small
\textit{Note.} Each panel reports RMSE by horizon for the nonlinear lag-quadratic DGP with $T=10,000$ to approximate asymptotic behavior. Rows correspond to persistence values $\rho \in \{0.5,0.7\}$. Columns within each panel correspond to nonlinearity strength $\gamma$. Estimator definitions and performance measures are identical to the baseline design.
\end{minipage}

\end{figure}

% ----------------------------------- Small Sample Deviation from Truth (Appendix)
\begin{figure}[htbp]
\centering
\caption{Small-Sample Deviation from the True IRF}
\label{fig:NON_LIN_SMALL_SAMPLE_DEV}

\begin{subfigure}[t]{0.9\textwidth}
    \centering
    \caption{\(\rho = 0.5\)}
    \includegraphics[width=\textwidth]{NON_LIN_LIMITED_DEVIATION_rho=0.5_T=250_B=5000.png}
    \label{fig:NON_LIN_SMALL_SAMPLE_DEV_rho05}
\end{subfigure}

\begin{subfigure}[t]{0.9\textwidth}
    \centering
    \caption{\(\rho = 0.7\)}
    \includegraphics[width=\textwidth]{NON_LIN_LIMITED_DEVIATION_rho=0.7_T=250_B=5000.png}
    \label{fig:NON_LIN_SMALL_SAMPLE_DEV_rho07}
\end{subfigure}

\begin{minipage}{0.9\textwidth}
\small
\textit{Note.} Deviation from truth, \(\hat\theta_h - \theta_h\), at \(T=250\) over \(B=5{,}000\) replications for the nonlinear lag-quadratic DGP, for VAR(4) (black) and LP(4) (crimson). Solid line is the mean deviation (finite-sample bias); shaded area is the central \(90\%\) Monte Carlo band (empirical 5th--95th percentile), not an SE band. Rows correspond to persistence \(\rho \in \{0.5,0.7\}\); columns to nonlinearity \(\gamma \in \{0,0.3,0.6\}\). The band widens with \(\gamma\) while the mean stays \(\gamma\)-stable.
\end{minipage}

\end{figure}


% ----------------------------------- Asymptotic Deviation from Truth (Appendix)
\begin{figure}[htbp]
\centering
\caption{Asymptotic Deviation from the True IRF}
\label{fig:NON_LIN_ASYMPTOTIC}

\begin{subfigure}[t]{0.9\textwidth}
    \centering
    \caption{\(\rho = 0.5\)}
    \includegraphics[width=\textwidth]{NON_LIN_ASYMPTOTIC_DEVIATION_rho=0.5_T=10000_B=1000.png}
    \label{fig:NON_LIN_ASYMPTOTIC_rho05}
\end{subfigure}

\begin{subfigure}[t]{0.9\textwidth}
    \centering
    \caption{\(\rho = 0.7\)}
    \includegraphics[width=\textwidth]{NON_LIN_ASYMPTOTIC_DEVIATION_rho=0.7_T=10000_B=1000.png}
    \label{fig:NON_LIN_ASYMPTOTIC_rho07}
\end{subfigure}

\begin{minipage}{0.9\textwidth}
\small
\textit{Note.} As above, at \(T=10{,}000\) over \(B=1{,}000\) replications. Both VAR(4) (black) and LP(4) (crimson) collapse to \(\approx 0\) at every horizon and \(\gamma\) --- consistent for the linear projection with no asymptotic bias floor --- so the residual misspecification is purely inferential. Solid line is the mean deviation; shaded area is the central \(90\%\) Monte Carlo band. The lag-quadratic is metastable, so divergent paths are dropped and the retained fraction is shown per panel. Rows correspond to \(\rho \in \{0.5,0.7\}\); columns to \(\gamma \in \{0,0.3,0.6\}\).
\end{minipage}

\end{figure}